\lecture{01}{Введение в распределённую оптимизацию}

\sourcevideo
    {https://www.youtube.com/playlist?list=PLOzRYVm0a65fhKDS8cYIC7YCuVGJ4FOJx}
    {Week 1: Lecture 1: Introduction to Optimization}

Распределённая оптимизация изучает задачи, в которых исходные данные,
вычисления и каналы связи распределены между несколькими агентами.
Каждый агент располагает лишь локальной информацией и обменивается
сообщениями с ограниченным числом соседей, но вся система должна
решить общую оптимизационную задачу.

\section{Архитектуры распределённого обучения}

Пусть имеется \(n\) агентов. Агент \(i\) хранит локальные данные
\(\Dataset_i\), локальную функцию \(f_i\) и переменную \(x_i\).
В централизованной постановке все функции доступны одному вычислителю,
который решает задачу
\[
    \min_{x\in\RR^d} F(x),
    \qquad
    F(x)\coloneqq\sum_{i=1}^{n}f_i(x).
\]

В распределённой постановке каждый агент работает только со своей
функцией. Эквивалентная задача записывается в виде
\[
    \begin{aligned}
        \min_{x_1,\ldots,x_n\in\RR^d}
        \quad&
        \sum_{i=1}^{n}f_i(x_i),
        \\
        \text{при условии}
        \quad&
        x_1=x_2=\cdots=x_n.
    \end{aligned}
\]
Равенство локальных переменных называется условием консенсуса.

\begin{definition}[Агент]
Агентом называется отдельный вычислительный узел, который хранит
локальные данные, выполняет локальные вычисления и обменивается
информацией с разрешёнными узлами сети.
\end{definition}

\begin{assumption}[Локальность информации]
Агент \(i\) непосредственно использует только собственную информацию
и сообщения агентов из множества соседей \(\neighbors{i}\).
\end{assumption}

Структура допустимых коммуникаций задаётся графом
\[
    \Graph=(\Vertices,\Edges),
    \qquad
    \Vertices=\{1,\ldots,n\}.
\]
В неориентированной модели наличие ребра
\((i,j)\in\Edges\) означает, что агенты \(i\) и \(j\)
могут обмениваться сообщениями.

Если вычисления выполняются локально, но все сообщения передаются
через единый сервер, архитектура является распределённой по
вычислениям, но централизованной по коммуникациям. Если и вычисления,
и обмен сообщениями происходят локально, архитектуру называют
децентрализованной.



Типичным примером распределённых вычислений служит обучение модели
предсказания текста на пользовательских устройствах. Устройство
хранит локальные данные и вычисляет градиент функции потерь, не
передавая сами данные центральному серверу.

Пусть \(w^k\in\RR^d\) --- параметры общей модели на итерации \(k\),
а \(L_i(w)\) --- функция потерь на данных устройства \(i\).
Сервер выбирает множество участвующих устройств \(S_k\), рассылает
им \(w^k\), получает локальные градиенты и выполняет обновление
\[
    w^{k+1}
    =
    w^k
    -
    \eta_k
    \frac{1}{\card{S_k}}
    \sum_{i\in S_k}
    \nabla L_i(w^k).
\]

Локальные вычисления распределены между устройствами, однако
агрегация остаётся централизованной. Частичное участие клиентов
позволяет не ожидать ответа от всех устройств на каждой итерации.
В простейшей модели предполагается, что локальные выборки имеют
близкие распределения.

Такая архитектура снижает необходимость передавать исходные данные,
но сохраняет центральный сервер как единую точку коммуникации и
потенциальную точку отказа.

\section{Экономическое распределение и необходимость координации}

Рассмотрим сеть из \(n\) электрогенераторов. Генератор \(i\)
производит мощность \(P_i\), а стоимость производства задаётся
квадратичной функцией
\[
    C_i(P_i)
    =
    a_iP_i^2+b_iP_i+c_i.
\]

\begin{problem}[Экономическое распределение мощности]
Требуется определить мощности генераторов из решения задачи
\[
    \begin{aligned}
        \min_{P_1,\ldots,P_n}
        \quad&
        \sum_{i=1}^{n}C_i(P_i),
        \\
        \text{при условиях}
        \quad&
        \sum_{i=1}^{n}P_i=P_{\mathrm{load}},
        \\
        &
        P_i^{\min}\le P_i\le P_i^{\max},
        \qquad i=1,\ldots,n.
    \end{aligned}
\]
Здесь \(P_{\mathrm{load}}\) --- суммарная требуемая мощность.
\end{problem}

Балансовое ограничение требует, чтобы суммарная генерация совпадала
с текущей нагрузкой. Ограничения
\(P_i^{\min}\le P_i\le P_i^{\max}\) описывают физические возможности
генераторов.

В централизованной постановке все коэффициенты \(a_i,b_i,c_i\)
известны одному оператору. В распределённой постановке генератор
знает только собственную функцию стоимости и может не раскрывать её
параметры другим участникам. Глобальное решение должно быть найдено
посредством локальных вычислений и обмена ограниченной информацией
между соседними генераторами.

Децентрализованная архитектура не требует центрального сервера.
Однако её отказоустойчивость зависит от топологии: отказ критической
вершины может нарушить связность сети. Коммуникационная нагрузка
отдельного агента определяется числом его соседей, а не общим числом
узлов сети.



Рассмотрим четыре локальные функции
\[
    f_i(x_i)=(x_i-i)^2,
    \qquad i=1,2,3,4.
\]
Без коммуникации каждый агент получает собственный минимум:
\[
    x_1^\star=1,
    \qquad
    x_2^\star=2,
    \qquad
    x_3^\star=3,
    \qquad
    x_4^\star=4.
\]

Если же агенты должны выбрать одно общее значение, возникает задача
\[
    \begin{aligned}
        \min_{x_1,x_2,x_3,x_4}
        \quad&
        \sum_{i=1}^{4}(x_i-i)^2,
        \\
        \text{при условии}
        \quad&
        x_1=x_2=x_3=x_4.
    \end{aligned}
\]
После подстановки общего значения \(x\) она принимает вид
\[
    \min_{x\in\RR}
    \sum_{i=1}^{4}(x-i)^2.
\]

\begin{lemma}[Минимум суммы квадратов]
Для чисел \(c_1,\ldots,c_n\in\RR\) функция
\[
    \Phi(x)=\sum_{i=1}^{n}(x-c_i)^2
\]
имеет единственный минимум в точке
\[
    x^\star=\frac{1}{n}\sum_{i=1}^{n}c_i.
\]
\end{lemma}

\begin{proof}
Имеем
\[
    \Phi'(x)
    =
    2\sum_{i=1}^{n}(x-c_i)
    =
    2n x-2\sum_{i=1}^{n}c_i.
\]
Условие \(\Phi'(x)=0\) даёт
\[
    x^\star
    =
    \frac{1}{n}\sum_{i=1}^{n}c_i.
\]
Поскольку \(\Phi''(x)=2n>0\), эта точка является единственным
минимумом.
\end{proof}

Для рассматриваемой задачи
\[
    x^\star
    =
    \frac{1+2+3+4}{4}
    =
    2.5.
\]
Следовательно, совокупность локальных минимумов не определяет
решение глобальной задачи. Для получения общего оптимума агентам
необходим обмен информацией.

\section{Консенсус и локальное усреднение}

Пусть каждый агент хранит скалярную оценку \(x_i^k\) на итерации
\(k\).

\begin{definition}[Консенсус]
Последовательность локальных оценок достигает консенсуса, если
существует \(c\in\RR\), для которого
\[
    \lim_{k\to\infty}x_i^k=c,
    \qquad
    i=1,\ldots,n.
\]
\end{definition}

\begin{definition}[Средний консенсус]
Достигается средний консенсус, если
\[
    \lim_{k\to\infty}x_i^k
    =
    \frac{1}{n}\sum_{j=1}^{n}x_j^0
\]
для каждого \(i\).
\end{definition}

Обычный консенсус требует лишь равенства предельных значений.
Средний консенсус дополнительно требует, чтобы общий предел совпадал
со средним начальных данных.



Пусть агент \(i\) получает оценки соседей и вычисляет
\[
    x_i^{k+1}
    =
    \sum_{j=1}^{n}a_{ij}x_j^k,
\]
где
\[
    a_{ij}=0
    \qquad
    \text{при }
    j\notin\neighbors{i}\cup\{i\}.
\]
Коэффициенты \(a_{ij}\) задают вес информации агента \(j\) при
обновлении агента \(i\).

Объединяя локальные оценки в вектор
\[
    x^k
    =
    \begin{pmatrix}
        x_1^k&\cdots&x_n^k
    \end{pmatrix}^{\mathsf T},
\]
получаем линейную итерацию
\[
    x^{k+1}=Ax^k,
    \qquad
    A=[a_{ij}].
\]

\begin{definition}[Матрица смешивания]
Матрицей смешивания называется неотрицательная матрица
\(A=[a_{ij}]\), определяющая обновление \(x^{k+1}=Ax^k\) и
удовлетворяющая условиям
\[
    a_{ij}=0
    \quad
    \text{при }
    j\notin\neighbors{i}\cup\{i\},
    \qquad
    \sum_{j=1}^{n}a_{ij}=1.
\]
\end{definition}

В лекции эта матрица называется матрицей смежности. Для
алгоритмического анализа точнее говорить о матрице смешивания,
поскольку её ненулевые элементы задают не только наличие связи,
но и вес передаваемой информации.

\begin{algorithm}[H]
\caption{Локальный алгоритм консенсуса}
\label{alg:lecture01-consensus}
\begin{algorithmic}[1]

\Require
граф \(\Graph=(\Vertices,\Edges)\),
матрица смешивания \(A=[a_{ij}]\),
начальные значения \(x_i^0\)

\For{\(k=0,1,2,\ldots\)}

    \ForAll{агентов \(i\in\Vertices\) параллельно}

        \State
        получить значения \(x_j^k\) от соседей
        \(j\in\neighbors{i}\)

        \State
        \(x_i^{k+1}
        \gets
        \displaystyle
        \sum_{j=1}^{n}a_{ij}x_j^k\)

    \EndFor

\EndFor

\Ensure
локальные последовательности \(\{x_i^k\}_{k\ge0}\)

\end{algorithmic}
\end{algorithm}

\section{Пример распределённого вычисления среднего}

Пусть четыре датчика имеют начальные измерения
\[
    x^0
    =
    \begin{pmatrix}
        25\\
        20\\
        24\\
        27
    \end{pmatrix}.
\]
Их среднее равно
\[
    \overline{x}^0
    =
    \frac{25+20+24+27}{4}
    =
    24.
\]

Рассмотрим матрицу смешивания
\[
    A_1
    =
    \begin{pmatrix}
        \frac12 & \frac12 & 0 & 0
        \\
        \frac14 & \frac14 & \frac14 & \frac14
        \\
        0 & \frac13 & \frac13 & \frac13
        \\
        0 & \frac13 & \frac13 & \frac13
    \end{pmatrix}.
\]
Каждая строка \(A_1\) суммируется в единицу, поэтому каждый агент
вычисляет взвешенное среднее доступных ему значений. Итерация
\[
    x^{k+1}=A_1x^k
\]
использует только локальные коммуникации.

В приведённом в лекции численном эксперименте после нескольких
итераций все компоненты сходятся приблизительно к одному значению:
\[
    x_i^k\longrightarrow 23.58.
\]
Следовательно, консенсус достигается, но его значение не равно
исходному среднему \(24\).

\begin{remark}
Нормировка строк матрицы смешивания обеспечивает корректность
локального усреднения, но сама по себе не гарантирует сохранение
глобального среднего.
\end{remark}



Определим среднее на итерации \(k\):
\[
    \overline{x}^k
    =
    \frac{1}{n}\one\trans x^k.
\]

\begin{lemma}[Инвариантность среднего]
Если матрица смешивания удовлетворяет условию
\[
    \one\trans A=\one\trans,
\]
то
\[
    \overline{x}^{k+1}=\overline{x}^k
\]
для всех \(k\).
\end{lemma}

\begin{proof}
Из равенства \(x^{k+1}=Ax^k\) следует
\[
    \overline{x}^{k+1}
    =
    \frac{1}{n}\one\trans x^{k+1}
    =
    \frac{1}{n}\one\trans Ax^k
    =
    \frac{1}{n}\one\trans x^k
    =
    \overline{x}^k.
\]
\end{proof}

\begin{theorem}[Консенсус при сохранении среднего]
Пусть среднее сохраняется и существует предел
\[
    x^k\longrightarrow c\one.
\]
Тогда
\[
    c=\overline{x}^0.
\]
Иными словами, достигнутый консенсус является средним консенсусом.
\end{theorem}

\begin{proof}
По инвариантности среднего
\[
    \overline{x}^k=\overline{x}^0.
\]
Переходя к пределу, получаем
\[
    \overline{x}^0
    =
    \lim_{k\to\infty}\frac{1}{n}\one\trans x^k
    =
    \frac{1}{n}\one\trans(c\one)
    =
    c.
\]
\end{proof}

В лекции для того же графа после изменения весов на достаточно
большой итерации численно получено
\[
    x_i^k\approx23.9989\approx24.
\]
Если новая матрица сохраняет среднее и обеспечивает консенсус, её
предельное значение равно точно \(24\). Следовательно, один и тот же
граф может приводить как к произвольному, так и к среднему консенсусу:
результат определяется не только топологией, но и весами обмена.

\section{Роль ускорения}

В прикладных задачах целевая точка может изменяться со временем.
Например, в задаче экономического распределения мощности величина
\(P_{\mathrm{load}}\) зависит от текущего потребления. Если алгоритм
сходится слишком медленно, найденное решение может соответствовать
уже устаревшему значению нагрузки.

Поэтому существенна не только сходимость метода, но и скорость
сходимости. Первая часть курса посвящена ускорению централизованных
методов оптимизации. Во второй части вводятся основы теории графов,
алгоритмы консенсуса и методы распределённого решения задач,
включая экономическое распределение мощности.


Распределённая оптимизация объединяет локальные вычисления и
ограниченный обмен информацией ради решения общей задачи.
Коммуникации необходимы, поскольку локальные минимумы в общем случае
не совпадают с глобальным решением. В алгоритмах консенсуса структура
графа определяет допустимые обмены, а матрица смешивания определяет
предельное значение. Обычный консенсус не гарантирует средний
консенсус; для последнего необходимо одновременно добиться
сходимости оценок и сохранить среднее начальных данных.